\section{Auswertung}
\label{sec:Auswertung}

\begin{figure}
  \centering
  \includegraphics{plot.pdf}
  \caption{Plot.}
  \label{fig:plot}
\end{figure}

\begin{table}
  \centering
  \caption{Eine Beispieltabelle mit Messdaten.}
  \label{tab:tabelle}
  \sisetup{table-format=1.1, per-mode=reciprocal}
  \begin{tblr}{
      colspec = {S[table-format=3.0] S[table-format=2.1] S},
      row{1} = {guard, mode=math},
      vline{4} = {2}{-}{text=\clap{$\pm$}},
    }
    \toprule
    U \mathbin{/} \unit{\volt} & I \mathbin{/} \unit{\micro\ampere} & \SetCell[c=2]{c} N \mathbin{/} \unit{\per\second} & \\
    \midrule
    360 & 0.1 & 98.3 & 0.9 \\
    400 & 0.2 & 99.8 & 1.0 \\
    420 & 0.2 & 99.1 & 0.9 \\
    \bottomrule
  \end{tblr}
\end{table}

Siehe \autoref{fig:plot} und \autoref{tab:tabelle}!

\subsection{Volumenbestimmung mit Maßstab}
\subsubsection*{Messdaten}

\subsubsection*{Herleitung der Formel}

Zur Bestimmung der Fromel haben wir zuerst den Körper in 2 Teile, einen oberen und einen unteren, aufgeteilt.\\
Es gilt also: Volumen des Körpers $=$ obere Teil $+$ untere Teil\\
oberer Teil:\\
Der obere Teil des Körpers basiert auf der Geometrie eines unvollständigen Zylinders. Wir berechnen also den vollständigen Zylinder und ziehen die Auskerungen ab.\\
gesamter Zylinder: V$=\pi\cdot r^{2} \cdot h = \pi \cdot \left(\frac{C_1}{2} \right)^2 \cdot B_6 $\\
Auskerbungen:\\
Schlitz A: $V = b \cdot l \cdot h = D_5 \cdot B_7 \cdot\left(\frac{D_1 - D_2 }{2} \right)$\\ 
Schlitz B: $V = b \cdot l \cdot h = D_4 \cdot B_8 \cdot\left(\frac{D_1 - D_2 }{2} \right) $\\ 
Zylinder oben: $V=\pi\cdot r^{2} \cdot h = \pi \cdot \left(\frac{D_2}{2} \right)^2 \cdot B_1$\\
Zylinder unten (klein): $V=\pi\cdot r^{2} \cdot h = \pi \cdot \left(\frac{A_7}{2} \right)^2 \cdot (B_6 - B_1)$\\
untere Teil:\\
Der untere Teil des Körpers basiert auf der Geometrie eines unvollständigen Quadars. Wir berechnen also den vollständigen Quadar und ziehen die Auskerungen bzw. Löcher ab.\\
gesamter Quadar: $V = a \cdot b \cdot c = A_1 \cdot A_9 \cdot B_5 $\\ 
Loch 1:
$V=\pi\cdot r_1^{2} \cdot h + \pi\cdot r_2^{2} \cdot h = \pi \cdot \left(\frac{D_3}{2} \right)^2 \cdot B_2 + \pi \cdot \left(\frac{A_2}{2} \right)^2 \cdot (B_5 - B_2) $\\
\subsubsection*{Formel für das Volumen}
